% Hafta 4 — Statik analiz ↔ dinamik analiz
% Derleme: py -3.12 tools/latex/derle.py docs/week-4/assets/h04-10-statik-dinamik.tex
\documentclass[tikz,border=8pt]{standalone}
\usepackage{cen429-cizim}
\begin{document}
\begin{tikzpicture}[node distance=9mm and 12mm]
  \node[baslik] (b) at (0,0) {Statik analiz $\leftrightarrow$ dinamik analiz};
  \node[altbaslik] at (b.south west) {biri diğerinin yerine geçmez --- ikisi birden kullanılır};

  \node[kutu=62mm, below=16mm of b.south west, anchor=north west] (sol)
    {\kb{STATİK (çalıştırmadan)}\\[2mm]
     tüm kod yollarına bakar\\
     hiç çalışmayan yolu da görür\\[3mm]
     yanlış pozitif ÇOK\\
     bağlamı bilmez};
  \node[iyikutu=62mm, right=22mm of sol] (sag)
    {\kb{DİNAMİK (çalıştırarak)}\\[2mm]
     yalnız yürütülen yol\\
     gerçek hatayı KESİN görür\\[3mm]
     yanlış pozitif AZ\\
     kapsam girdiye bağlı};

  \node[fit=(sol)(sag), inner sep=0] (ikifit) {};
  \coordinate (orta) at ($(sol.east)!0.5!(sag.west)$);
  \draw[cizgi, line width=1pt] (ikifit.north -| orta) -- (ikifit.south -| orta);

  \node[dip, text width=170mm, below=10mm of ikifit.south, anchor=north]
    {Fuzzing üçüncü ayaktır: dinamik analizin göreceği YENİ yolları üretir.};
\end{tikzpicture}
\end{document}
